\documentclass[journal]{IEEEtran}

% Text encoding and fonts
\usepackage[utf8]{inputenc}
\usepackage[T1]{fontenc}
\usepackage{microtype}

% Math packages
\usepackage{amsmath}
\usepackage{amsfonts}

% Graphics and figures
\usepackage{graphicx}
\graphicspath{{figs/}}
 
% Tables
\usepackage{booktabs}
\usepackage{multirow}

% Algorithms
\usepackage{algorithm}
\usepackage[noEnd=false,indLines=true]{algpseudocodex}

% Spacing and formatting
% \usepackage{setspace}
\allowdisplaybreaks[2]

% Hyperlinks and URLs
\usepackage{hyperref}
\usepackage{url}

% ========================================
% DOCUMENT METADATA
% ========================================
\begin{document}

% Title
\title{Accelerating Underground Pumped Hydro Energy Storage Scheduling with Decision-Focused Learning}

% Authors
\author{Honghui~Zheng,
        Pietro~Favaro,
        Yury~Dvorkin, %~\IEEEmembership{Member,~IEEE,}
        and~Ján~Drgoňa, %~\IEEEmembership{Member,~IEEE}
\thanks{H. Zheng, Y. Dvorkin, and J. Drgoňa are with the Department of Civil and Systems Engineering, Johns Hopkins University, Baltimore, MD 21218, USA (e-mail: \{hzheng39, ydvorki1, jdrgona1\}@jh.edu).}
\thanks{P. Favaro is with the Power Systems and Markets Research Group, University of Mons, 7000 Mons, Belgium (e-mail: pietro.favaro@umons.ac.be).}
\thanks{This work is under review.}}

% The paper headers
\markboth{Preprint}%
{Running Title for Header}

\maketitle

% ========================================
% ABSTRACT
% ========================================
\begin{abstract}
Underground pumped hydro energy storage (UPHES) systems play a critical role in grid-scale energy storage for renewable integration, yet optimal day-ahead scheduling remains computationally prohibitive due to nonlinear turbine performance characteristics and discrete operational modes. This paper presents a decision-focused learning (DFL) framework that addresses the computational-accuracy trade-off in UPHES day-ahead scheduling. The proposed methodology employs neural networks to predict penalty weights that guide recursive linearization, transforming the intractable MINLP into a sequence of convex quadratic programs trained end-to-end via differentiable optimization layers. Case studies across 19 representative Belgian electricity market scenarios demonstrate that the DFL framework effectively navigates the trade-off between solution quality and computation time. As a refinement tool, the framework improves profit by 1.1\% over piecewise MIQP baselines. Alternatively, as a real-time scheduler initialized with linear approximations, it achieves a 300-fold speedup (3.87~s vs 1205.79~s) while maintaining profitability within 3.6\% of the piecewise MIQP benchmark. Thus, the presented DFL framework enables flexible prioritization between profit maximization and real-time responsiveness. 
\end{abstract}

% Keywords
\begin{IEEEkeywords}
Data-driven optimization, Decision-focused learning, Pumped hydro energy storage \end{IEEEkeywords}

\section*{Nomenclature}
\addcontentsline{toc}{section}{Nomenclature}

\subsection*{Sets and Indices}
\begin{IEEEdescription}[\IEEEusemathlabelsep\IEEEsetlabelwidth{$\mathcal{M}_{\mathrm{act}}$}]
\item[$\mathcal{T}$] Set of time periods in the planning horizon.
\item[$t$] Time index representing hourly intervals, $t \in \mathcal{T}$.
\item[$\mathcal{M}$] Set of operating modes, $\mathcal{M} = \{I, T, P\}$ (idle, turbine, pump).
\item[$m$] Operating-mode index, $m \in \mathcal{M}$.
\item[$\mathcal{I}_h$] Set of head discretization points for SOS2.
\item[$\mathcal{I}_p^m$] Set of power discretization points in SOS2 for $m$.
\item[$i, j$] Indices for discretization points in SOS2.
\end{IEEEdescription}

\subsection*{Decision Variables}
\begin{IEEEdescription}[\IEEEusemathlabelsep\IEEEsetlabelwidth{$\Omega_{t,i,j}^m$}]
\item[$p_t^m$] Power output (turbine) or consumption (pump) at time $t$ for mode $m \in \{T, P\}$ [MW].
\item[$q_t^m$] Water flow rate at time $t$ for mode $m$ [m$^3$/s].
\item[$h_t$] Hydraulic head at time $t$ [m].
\item[$v_{\text{low},t}$] Lower reservoir volume at time $t$ [m$^3$].
\item[$z_t^m$] Binary status variable for mode $m \in \mathcal{M}$ at time $t$.
\item[$\Omega_{t,i}^{\text{vh}}$] SOS2 interpolation weight for volume-head relationship at time $t$, discretization point $i$.
\item[$\Omega_{t,i,j}^m$] SOS2 interpolation weight for UPC of mode $m$ at time $t$, grid point $(i,j)$.
\end{IEEEdescription}

\subsection*{System Parameters}
\begin{IEEEdescription}[\IEEEusemathlabelsep\IEEEsetlabelwidth{$h_\text{min}, h_\text{max}$}]
\item[$\lambda_t^{\text{DA}}$] Day-ahead electricity price at time $t$ [€/MWh].
\item[$C_{\text{op}}$] Operational cost coefficient [€/MW$^2$].
\item[$\Delta t$] Time step duration, 3600 s.
\item[$v_{\text{low}}^{\text{init}}$] Initial lower reservoir volume [m$^3$].
\item[$v_{\text{low}}^{\text{target}}$] Target lower reservoir volume [m$^3$].
\item[$v_{\text{total}}$] Total combined capacity of reservoirs [m$^3$].
\item[$h_\text{min}, h_\text{max}$] Minimum and maximum hydraulic head [m].
\item[$h^{\text{init}}$] Initial hydraulic head [m].
\end{IEEEdescription}

\subsection*{Nonlinear Functions and Mappings}
\begin{IEEEdescription}[\IEEEusemathlabelsep\IEEEsetlabelwidth{$f_m^{\text{UPC}}(p, h)$}]
\item[$f_m^{\text{UPC}}(p, h)$] Unit performance curve for mode $m$ [m$^3$/s].
\item[$p_{\min}^m(h)$] Minimum power for mode $m$ [MW].
\item[$p_{\max}^m(h)$] Maximum power for mode $m$ [MW].
\item[$f_{\text{up}}^{\text{vol}}(h)$] Upper reservoir volume function [m$^3$].
\item[$f_{\text{low}}^{\text{vol}}(h)$] Lower reservoir volume function [m$^3$].
\item[$c_{ab}^{(m)}$] Polynomial coefficients for UPC of mode $m$.
\item[$d$] Total polynomial degree for UPC.
\end{IEEEdescription}

\subsection*{Linearization Coefficients}
\begin{IEEEdescription}[\IEEEusemathlabelsep\IEEEsetlabelwidth{$\boldsymbol{\beta}_m^{\max}$}]
\item[$\boldsymbol{\alpha}_m$] Coefficient vector for global linear approximation of UPC for mode $m$ [m$^3$/s per (MW, m)].
\item[$\boldsymbol{\beta}_m^{\min}$] Coefficient vector for linearization of minimum power boundary, mode $m$ [MW per m].
\item[$\boldsymbol{\beta}_m^{\max}$] Coefficient vector for linearization of maximum power boundary, mode $m$ [MW per m].
\item[$\boldsymbol{\delta}$] Coefficient vector for global linear approximation of volume-head relationship [m per m$^3$].
\item[$\boldsymbol{\xi}_{q,t}^{(k)}$] Local linearization coefficient vector for UPC at time $t$, iteration $k$.
\item[$\boldsymbol{\xi}_{v,t}^{(k)}$] Local linearization coefficient vector for volume-head at time $t$, iteration $k$.
\end{IEEEdescription}

\subsection*{Sample Points (Piecewise Interpolation)}
\begin{IEEEdescription}[\IEEEusemathlabelsep\IEEEsetlabelwidth{$q_{i,j}^{m,\text{sample}}$}]
\item[$h_i^{\text{sample}}$] Sampled head value at point $i$ [m].
\item[$v_{\text{low},i}^{\text{sample}}$] Sampled lower reservoir volume at point $i$ [m$^3$].
\item[$p_{i,j}^{m,\text{sample}}$] Sampled power value at point $(i,j)$ for $m$ [MW].
\item[$q_{i,j}^{m,\text{sample}}$] Sampled flow rate at point $(i,j)$ for $m$ [m$^3$/s].
\end{IEEEdescription}

\subsection*{Decision-Focused Learning Parameters}
\begin{IEEEdescription}[\IEEEusemathlabelsep\IEEEsetlabelwidth{$\mathbf{w}_p, \mathbf{w}_q, \mathbf{w}_h$}]
\item[$\theta$] Neural network parameters.
\item[$\mathbf{w}_p, \mathbf{w}_q, \mathbf{w}_h$] Penalty weights for power, flow, and head.
\item[$K$] Total number of recursive linearizations.
\item[$\gamma > 1$] Penalty growth factor.
\end{IEEEdescription}

\subsection*{Performance Metrics}
\begin{IEEEdescription}[\IEEEusemathlabelsep\IEEEsetlabelwidth{$\text{Vol}$}]
\item[$\Pi$] Ex-post profit evaluated using simulated trajectories under true nonlinear dynamics [EUR].
\item[$\text{SI}$] System imbalance penalty for power schedule [EUR].
\item[$\text{Vol}$] Penalty for exceeding target volume [EUR].
\item[$\mathcal{L}(\theta)$] Training loss function (negative ex-post profit).
\end{IEEEdescription}

\subsection*{Solution Notation Convention}
\begin{IEEEdescription}[\IEEEusemathlabelsep\IEEEsetlabelwidth{$\hat{\mathbf{x}}^{(k)}$}]
\item[$\bar{\mathbf{x}}$] Initial noisy solution from baseline MIQP.
\item[$\hat{\mathbf{x}}^{(k)}$] Optimized solution at iteration $k$.
\item[$\tilde{\mathbf{x}}$] Simulated solution under true nonlinear dynamics.
\end{IEEEdescription}

% ========================================
% SECTION 1: INTRODUCTION
% ========================================
\section{Introduction} 
\label{sec:introduction}

\IEEEPARstart{G}{rid}-scale energy storage is critical to compensate for the temporal variance of weather-dependent generation, such as wind, solar, and loads~\cite{golshani2018coordination}. Despite the recent proliferation of battery storage systems~\cite{wang2016review}, pumped hydro energy storage (PHES) is anticipated to remain the dominant form of grid-scale energy storage globally, particularly for durations of four hours and more. As a proven technology with decades of operational history~\cite{REHMAN2015586}, PHES accounts for the vast majority of installed large-scale storage capacity, representing about 99\% of the world’s electricity storage capacity and around 3\% of the total installed electricity generation capacity globally~\cite{rastler2010electricity}, far surpassing -- at least for the present moment -- the contribution of battery energy storage systems~\cite{blakers2021review}. Both the operational maturity and deployment scale have established PHES as the backbone of grid-scale energy storage, providing reliable bulk energy shifting and ancillary services~\cite{hunt2020existing}.

Although conventional PHES installations require specific landscapes with significant elevation differences between upper and lower reservoirs, Underground Pumped Hydro Energy Storage (UPHES) is suited for regions with limited topographical relief or environmental constraints that preclude natural surface reservoirs or their low-impact development~\cite{menendez2019pumped}. For example, by repurposing abandoned mines~\cite{winde2017exploring}, and aquifers~\cite{martin2007aquifer} as lower reservoirs, UPHES circumvents such geographical limitations and reduces environmental impacts and land-use requirements~\cite{yang2016pumped}. UPHES also has the potential to transform decommissioned mining facilities into storage assets~\cite{pickard2011history}, which in turn may support local communities~\cite{TASSI2024114467, COLAS2023109153}.

Operating the UPHES system efficiently is a substantial computational challenge due to nonlinear physical dynamics and relatively large round-trip inefficiencies~\cite{emrani2022assessment}, discrete operational turbine and pumping modes~\cite{jia2013mixed}, and optimization under uncertain electricity market conditions~\cite{benini2002day, 7353197}. The original computational challenge stems from the Unit performance curves (UPCs) which relate hydraulic head (elevation difference between upper and lower reservoirs), power output, and water flow via a complex nonlinear and non-convex mapping~\cite{finardi2016assessing}. Furthermore, underground reservoir geometries introduce discontinuous head variations as water levels transition between horizontal galleries, creating stepwise changes in available head and turbine efficiency~\cite{toubeau2019non}. The system must transition among discrete operational states (pumping, generating, and idle) subject to mechanical constraints and market incentives that vary hourly or sub-hourly in electricity markets~\cite{cheng2022hierarchical}.

Existing optimization approaches for PHES scheduling predominantly employ mixed-integer programming (MIP) formulations~\cite{zhang2025long, jia2013mixed} that approximate nonlinear relationships through piecewise linearization~\cite{toubeau2019non, toubeau2019chance}. While these formulations leverage powerful solvers and provide optimality guarantees, they face fundamental trade-offs between approximation fidelity and computational tractability~\cite{andrade2022integer}. Piecewise approximations with fine discretization lead to improved solutions at the expense of substantial computational burdens that may hinder real-time application~\cite{vielma2010mixed}. Recent advances in neural network (NN)-constrained optimization~\cite{favaro2024neural} have demonstrated that trained NNs can accurately model nonconvex UPCs, but introduced additional integer variables proportional to the number of neurons in the network architecture. Despite significant progress in MIP modeling, these approximations struggle to capture nonlinearities, often yielding infeasible schedules. Alternative approaches using nonlinear optimization coupled with high-fidelity physical simulators~\cite{favaro_multi-fidelity_2024} can discover superior solutions by directly optimizing ex-post operational performance, but require substantial computational resources that preclude real-time deployment. This highlights the need for optimization frameworks that leverage learned problem structure to maintain both solution accuracy and computational efficiency.

Decision-focused learning (DFL) offers an approach to this challenge by directly training models to minimize downstream decision costs rather than nominal prediction errors~\cite{mandi2024decision, elmachtoub2022smart}. Traditional two-stage approaches train predictive models to minimize statistical errors under the assumption that estimation accuracy translates to optimal decisions~\cite{bertsimas2018predictive}. However, this assumption fails in complex energy systems where asymmetric cost structures and nonlinearities create non-monotonic relationships between prediction accuracy and decision quality~\cite{zhang2025decision}. Power system case studies~\cite{zhang2025decision} show certain estimation errors may incur negligible operational penalties while others critically compromise solution feasibility, yet standard two-stage models treat estimation quality uniformly~\cite{demirovic2019investigation}. 

DFL addresses this limitation of the predictive models by embedding the optimization model within the training pipeline through differentiable optimization layers~\cite{cvxpylayers2019}, enabling end-to-end learning with loss functions derived from operational objectives or constraints rather than statistical fit. Recent power system applications demonstrate DFL's operational advantages, including wind power forecasting integrated with energy scheduling~\cite{wahdany2023more}, electricity price forecasting~\cite{tschora2023forecasting}, and demand response~\cite{favaro_decision-focused_2025}. Benchmark comparison reveals that DFL reduces operational costs and decision regret by 15\%--40\% compared to two-stage approaches~\cite{zhang2025decision, mandi2024decision}. For UPHES systems, the combination of UPC, discrete operational modes, and complex reservoir behavior creates conditions where DFL's direct optimization of decision quality offers substantial advantages over conventional forecasting.
This paper seeks to enhance the operational utility of UPHES by implementing day-ahead optimal scheduling using an end-to-end DFL approach. 

\noindent The contributions of this work are as follows:

\begin{enumerate}
    \item An end-to-end DFL framework is proposed for optimal day-ahead UPHES scheduling, utilizing differentiable optimization. A recursive local linearization technique refines an initial scheduling solution obtained from a conventional mixed-integer optimization method or a historical operation schedule.
    \item A neural penalty weight predictor is introduced to ensure solution feasibility by mapping optimization parameters to the trust region size of a given local linearization.
    \item A differentiable physics simulator is developed to directly evaluate the ex-post operational profit as training loss, ensuring solution quality over training.
    \item Case studies demonstrated that the DFL framework effectively navigates the accuracy-computation trade-off: as a refinement tool, it improves profit by 1.1\%; alternatively, as a real-time scheduler, it achieves a 300-fold speedup while maintaining profitability within 3.6\% of the state-of-the-art piecewise linear MIQP  baseline.
    \item We open-source our code implementation of the proposed DFL framework to promote reproducibility and adoption\footnote{\url{https://github.com/SOLARIS-JHU/DFL-UPHES}}.
\end{enumerate}

\section{Problem Formulation}
\label{sec:problem}
This section presents the mathematical formulation of the UPHES scheduling problem in \eqref{eq:uphes_scheduling_problem} in the context of the day-ahead electricity market. This optimization is formulated as a Mixed-Integer Nonlinear Programming (MINLP) problem to account for nonlinear and discrete characteristics of the UPC and volume-head relationships in the lower reservoir.

Fig.~\ref{fig:uphes_system} illustrates the system configuration. In turbine mode, water descends to generate electricity; in pump mode, electrical power lifts water to the upper reservoir. 

\begin{figure}[t!]
    \centering
    \includegraphics[width=\columnwidth]{figs/UPHES.pdf}
    \caption{Schematic of the UPHES system showing the reservoirs, reversible Francis turbine, and gross head $h_t$ bounded by $h_{\min}$ and $h_{\max}$.}
    \label{fig:uphes_system}
\end{figure}

\begin{subequations}
\label{eq:uphes_scheduling_problem}
\begin{gather}
\underset{p_t, z_t^I, z_t^T, z_t^P}{\text{maximize}} \quad \sum_{t=1}^{24} \Delta t \left( p_t \cdot \lambda^{\text{DA}}_t - C_{\text{op}} \cdot p_t^2 \right) \label{eq:objective_complete} \\[2pt]
\text{s.t.} \qquad \sum_{m \in \mathcal{M}} z_t^m = 1, \quad z_t^m \in \{0,1\}, \quad \forall t \in \mathcal{T} \label{eq:mode_sum_binary}\\
p_t = p_t^T + p_t^P,\quad q_t = q_t^T + q_t^P,\quad \forall t \in \mathcal{T} \label{eq:agg_power_flow}\\
 m = I, \ z_t^m = 1 \ \Rightarrow\  p_t = 0,\ q_t = 0, \quad \forall t \in \mathcal{T} \label{eq:idle_indicator}\\
(p_t^m, q_t^m) = 
\begin{cases}
\big(p \in [p_{\min}^m(h_t), p_{\max}^m(h_t)],\ f_m^{\text{UPC}}(p_t^m, h_t)\big), \\
\phantom{(0, 0),} \hspace{0.7em} \text{if } z_t^m = 1, \forall t \in \mathcal{T},\ \forall m \in \{T, P\} \\
(0, 0), \hspace{0.5em} \text{if } z_t^m = 0, \forall t \in \mathcal{T},\ \forall m \in \{T, P\}
\end{cases} 
\label{eq:upc_constraints}\\
v_{\text{low},1} = v_{\text{low}}^{\text{init}} + \Delta t \cdot q_1 \label{eq:volume_initial} \\[6pt]
v_{\text{low},t} = v_{\text{low},t-1} + \Delta t \cdot q_t, \quad \forall t \in \{2,3,\ldots,24\} \label{eq:volume_dynamics} \\
v_{\text{low},t} = f^{\text{vol}}(h_t), \quad \forall t \in \mathcal{T} \label{eq:volume_head} \\
h_{\min} \le h_t \le h_{\max}, \quad \forall t \in \mathcal{T} \label{eq:head_bounds} \\
v_{\text{low},24} \leq v_{\text{low}}^{\text{target}} \label{eq:Target}
\end{gather}
\end{subequations}

Here, $t \in \mathcal{T}=\{1,\dots,24\}$ indexes hourly intervals, and $m \in \mathcal{M}=\{I, T, P\}$ denotes the operating mode (Idle, Turbine, Pump), with binary indicators $z_t^m$. Continuous variables $p_t$, $q_t$, $h_t$, and $v_{\text{low},t}$ represent net power [MW], flow [m$^3$/s], hydraulic head [m], and lower reservoir volume [m$^3$], while $\lambda_t^{\text{DA}}$ and $C_{\text{op}}$ define the market price and cost coefficient.

The objective function \eqref{eq:objective_complete} maximizes the operator's profit over a 24-hour horizon. The operator participates as a price taker in the day-ahead energy market, which clears hourly electricity prices $\lambda_t^{\text{DA}}$ for the following day. The revenue term $p_t \cdot \lambda^{\text{DA}}_t$ captures the electricity production value (positive for turbine operation, negative for pumping). The quadratic operational cost term $C_{\text{op}} \cdot p_t^2$ reflects the nonlinear operational expenses, including mechanical wear, friction losses, and efficiency degradation at higher loading~\cite{yin2023predictive, aburub2019use}. The coefficient $C_{\text{op}}$ effectively aggregates these degradation mechanisms into a scalar operational cost, while $\Delta t$ is the hourly time step.

\subsection{Operational Constraints for Francis Turbine}

Fig.~\ref{fig:UPC_poly_fit_2D} illustrates the unit performance curves (UPCs) obtained from laboratory measurements on a reduced-scale reversible Francis pump–turbine reported by Mercier et~al.~\cite{mercier2017provision} within the SmartWater project~\cite{smartwater_multitel_2022}. The operating space consists of two disjoint manifolds (turbine and pump) separated by an idle transition at zero flow. This discontinuous, non-convex geometry reflects distinct hydraulic regimes: turbine mode extracts energy from descending flow, pump mode lifts water against gravity, and idle mode maintains zero flow during transitions. The mapping between head, power, and flow is consequently non-linear, non-convex, and discontinuous across these operational domains. 
Making an accurate representation of these performance characteristics is essential for optimal scheduling. Francis pump–turbines are employed for their rapid load regulation~\cite{iliev2019variable}, suitability for high-head applications~\cite{wang2020transient}, and established use in utility-scale energy storage~\cite{koritarov2022review}.

\begin{figure}[t!]
    \centering
    \includegraphics[width=\linewidth]{figs/francis_turbine_polynomial_combined_contours.pdf}
    \caption{Polynomial fit of the unit performance curves: (left) Pump mode, (right) Turbine mode.}
    \label{fig:UPC_poly_fit_2D}
\end{figure}

To embed the UPCs into an optimization framework, we approximate the experimental data with a bivariate polynomial that maps head and power to discharge. The polynomial degree $d=5$ was selected through evaluation of approximation accuracy versus model complexity. We obtain the mode-dependent surrogate
\begin{equation}
\label{eq:upc-poly}
q_t^{m} = f_m^{\mathrm{UPC}}(p_t,h_t)
=
\sum_{a=0}^{d}\sum_{b=0}^{d-a} c^{m}_{ab}\,p_t^{\,a}\,h_t^{\,b},
\quad m\in\{T,P\},
\end{equation}
where the coefficients $c^{m}_{ab}$ are estimated by least squares. Because turbine and pump samples occupy distinct regions of the $(p_t,h_t)$ plane, \eqref{eq:upc-poly} is fitted separately for each mode. The polynomial approximation achieves high fidelity with $R^2 \geq 0.99988$ for both turbine and pump modes. The projected envelopes are shown in Fig.~\ref{fig:UPC_poly_fit_2D}.

\subsubsection{Mode Selection Constraints}
Constraint~\eqref{eq:mode_sum_binary} ensures mutual exclusivity among operational modes through binary indicators $z_t^m$, at each time $t$. These indicators govern the aggregated power $p_t$ and flow $q_t$ variables defined in~\eqref{eq:agg_power_flow}.

\subsubsection{Idle Mode Constraints}
When idle mode is engaged ($z_t^I=1$), constraint~\eqref{eq:idle_indicator} enforces $p_t=0$ and $q_t=0$, eliminating power exchange and hydraulic flow. Through~\eqref{eq:agg_power_flow}, this condition propagates to both mode-specific variables, while~\eqref{eq:upc_constraints} ensures $(p_t^T, q_t^T) = (0,0)$ and $(p_t^P, q_t^P) = (0,0)$ when the respective mode indicators are inactive.


\subsubsection{Active Mode Constraints}
For active modes $m \in \{T, P\}$ with $z_t^m=1$, constraint~\eqref{eq:upc_constraints} bounds power within $[p_{\min}^m(h_t), p_{\max}^m(h_t)]$, and enforces flow rate $q_t^m = f_m^{\text{UPC}}(p_t^m, h_t)$ via ~\eqref{eq:upc-poly}. When $z_t^m=0$, both $p_t^m$ and $q_t^m$ are set to zero, decoupling the inactive mode from system dynamics. The head-dependent power boundaries ensure operation remains within the hydraulically admissible envelope, as operating outside this range is hazardous: excessive flow can trigger damaging cavitation (rapid implosion of gaseous cavities) in the Francis turbine, whereas insufficient flow can lead to self-excited vibrations and severe erosion~\cite{pannatier2010optimisation}.

\subsection{Volume Dynamics Constraints}
\label{sec:volume_constraints}

The lower reservoir volume evolves according to volume conservation, where the rate of volume change equals the net water flow: $\frac{dv_{\text{low}}}{dt}=q(t)$. Constraints~\eqref{eq:volume_dynamics} and~\eqref{eq:volume_initial} discretize this continuous balance using forward Euler integration with time step $\Delta t = 3600$~s, linking hourly flow decisions to cumulative storage levels throughout the scheduling horizon.

The nonlinear volume–head coupling arises from the modeling assumptions on reservoir geometries that deliberately introduce nonlinearity and nonconvexity into the optimization problem. These geometric configurations capture the computational challenges inherent in real UPHES systems while maintaining analytical tractability. Following the frustum-based storage-elevation modeling approach~\cite{liu2022low}, the upper reservoir is represented as a frustum with average slope $m$, yielding the cubic volume-head relationship
\begin{align}
v_{\text{up}}(h_{\text{up}}) = \pi r_{\text{base}}^2 h_{\text{up}} + \pi m r_{\text{base}} h_{\text{up}}^2 + \frac{\pi m^2}{3} h_{\text{up}}^3. \label{eq:v_up_cubic}
\end{align}
For the lower reservoir, various underground configurations have been proposed for UPHES systems, including room-and-pillar layouts and interconnected cavern networks~\cite{allen1984underground}. Following established modeling approaches, we represent the lower reservoir as $n$ spherical mine pits, yielding 
\begin{align}
v_{\text{low}}(h_{\text{low}}) = n\pi R h_{\text{low}}^2 - \frac{n\pi}{3} h_{\text{low}}^3.\label{eq:v_low_cubic}
\end{align}
The gross head is given as 
\begin{align}
h = f_{\text{up}}^{\text{vol}\,-1}(v_{\text{up}}) - f_{\text{low}}^{\text{vol}\,-1}(v_{\text{low}}), \label{eq:head_vol_relation}
\end{align}
subject to $v_{\text{up}} + v_{\text{low}} = v_{\text{total}}$, where $f_{\text{up}}^{\text{vol}\,-1}$ and $f_{\text{low}}^{\text{vol}\,-1}$ denote the inverse mappings obtained by solving the cubic equations \eqref{eq:v_up_cubic} and (\ref{eq:v_low_cubic}), respectively.

The hydraulic head is constrained within its feasible operating range as enforced by~\eqref{eq:head_bounds}. This condition ensures that the computed head levels remain physically realistic and consistent with the reservoir geometry throughout the scheduling horizon.

Target constraint \eqref{eq:Target} ensures that the final lower reservoir volume $v_{\text{low},t=24}$ does not exceed the target $v_{\text{low}}^{\text{target}}$, preserving long-term water balance and operational consistency.

\section{Preliminaries}
\label{sec:preliminaries}

The MINLP described in Section~II poses computational challenges due to polynomial UPC mappings~\eqref{eq:upc-poly}, cubic volume-head relations~\eqref{eq:v_up_cubic} and~\eqref{eq:v_low_cubic}, and head-dependent power boundaries. This section presents two tractable piecewise reformulations of these nonlinear relationships that retain the binary mode variables from~\eqref{eq:mode_sum_binary}. These formulations serve as baselines for evaluating the DFL methodology in Section~\ref{sec:methodology}.

\subsection{Global Linearization Approach}

Global linearization approximates each nonlinear relationship with a single affine function fitted over the entire operational domain, trading accuracy for computational efficiency. For the UPC mappings in~\eqref{eq:upc-poly}, we replace the polynomial with
\begin{align}
q_t^m = 
\begin{bmatrix} 
p_t^m & h_t & 1 
\end{bmatrix}
\boldsymbol{\alpha}_m,
\quad \forall t \in \mathcal{T},\ m \in \{T,P\}
\label{eq:upc_linear}
\end{align}
where $\boldsymbol{\alpha}_m \in \mathbb{R}^3$ is obtained by least-squares fitting on experimental data. The power boundaries are linearized as
\begin{align}
p_{\min}^m(h_t) &= 
\begin{bmatrix} 
h_t & 1 
\end{bmatrix}
\boldsymbol{\beta}_m^{\min}, 
\quad
p_{\max}^m(h_t) = 
\begin{bmatrix} 
h_t & 1 
\end{bmatrix}
\boldsymbol{\beta}_m^{\max}, \label{eq:pminmax_linear}
\end{align}
for $m \in \{T,P\}$. The volume-head relationship~\eqref{eq:volume_head} becomes
\begin{align}
h_t = 
\begin{bmatrix} 
v_{\text{low},t} & 1 
\end{bmatrix}
\boldsymbol{\delta}, 
\quad \forall t \in \mathcal{T}
\label{eq:vh_linear}
\end{align}
where $\boldsymbol{\delta} \in \mathbb{R}^2$. Idle-mode constraints~\eqref{eq:idle_indicator} are enforced via big-$M$ formulations. While computationally efficient, the single-plane approximation introduces systematic errors in high-curvature regions, motivating the adaptive local linearization strategy developed in Section~\ref{sec:methodology}.

\subsection{Piecewise Bilinear Approximation with SOS2 Constraints}

This second method enhances the fidelity of the UPC approximation while retaining a tractable formulation. We discretize the function space on the variable grid and enforce interpolation only between adjacent grid vertices. The resulting interpolant is a bilinear hyperbolic paraboloid. A Special Ordered Set of type 2 (SOS2) constraint~\cite{beale1970special} is employed to ensure that the convex‑combination weights are non‑zero only for two neighboring vertices in each cell~\cite{beale1976global}. This approach is grounded in modeling of non‑separable multivariate functions in a mixed-integer manner~\cite{huchette2019combinatorial,barmann2023piecewise}.

For the volume-head, SOS2 variables $\Omega_{t,i}^{\text{vh}} \in [0,1]$ satisfy
\begin{equation}
\sum_{i \in \mathcal{I}_h} \Omega_{t,i}^{\text{vh}} = 1, \quad
(h_t, v_{\text{low},t}) = \sum_{i \in \mathcal{I}_h} \Omega_{t,i}^{\text{vh}} (h_i^{\text{sample}}, v_{\text{low},i}^{\text{sample}}),
\label{eq:sos2_vh}
\end{equation}
where at most two adjacent $\Omega_{t,i}^{\text{vh}}$ are nonzero. For UPC mappings, auxiliary variables $\Omega_{t,i,j}^m \in [0,1]$ satisfy
\begin{equation}
\sum_{j \in \mathcal{I}_p^m} \Omega_{t,i,j}^m = z_t^m \cdot \Omega_{t,i}^{\text{vh}}, \quad
\sum_{i \in \mathcal{I}_h}\sum_{j \in \mathcal{I}_p^m} \Omega_{t,i,j}^m = z_t^m,
\label{eq:sos2_mode}
\end{equation}
with power and flow reconstructed as
\begin{equation}
(p_t^m, q_t^m) = \sum_{i \in \mathcal{I}_h}\sum_{j \in \mathcal{I}_p^m} \Omega_{t,i,j}^m (p_{i,j}^{m,{\text{sample}}}, q_{i,j}^{m,{\text{sample}}}).
\label{eq:sos2_power_flow}
\end{equation}
This formulation improves accuracy over global linearization at the cost of $O(N_h \cdot N_p^T + N_h \cdot N_p^P)$ auxiliary variables, where $N_h$ denotes the number of head discretization points and $N_p^m$ represents the number of power discretization points for mode $m \in \{T, P\}$. The computational complexity scales linearly with grid resolution, requiring $24(N_h \cdot N_p^T + N_h \cdot N_p^P)$ SOS2 interpolation weights $\Omega_{t,i,j}^m$ over the planning horizon.

\section{Methodology}
\label{sec:methodology}

This section describes a DFL framework for computationally efficient and accurate UPHES day-ahead scheduling. As illustrated in Fig.~\ref{fig:methodology_framework} and detailed in Algorithm~\ref{alg:dfl_scheduling}, the methodology integrates four components: (\textit{i}) a neural penalty predictor (Section~\ref{subsec:penalty_prediction}) that generates adaptive regularization weights to establish trust regions ensuring linearization validity, (\textit{ii}) a local linearization layer (Section~\ref{subsec:local_linearization}) that constructs first-order Taylor approximations converting nonlinear constraints into convex form, (\textit{iii}) a differentiable convex optimizer (Section~\ref{subsec:diff_optimizer}) that iteratively refines schedules through penalty-guided quadratic programming while preserving mode consistency, and (\textit{iv}) a physical simulator (Section~\ref{subsec:simulation}) that enforces true nonlinear dynamics to validate feasibility and quantify the economic performance used as training signal. 
This end-to-end differentiable DFL architecture enables optimization of economic performance across a range of price scenarios and can be used as a computationally tractable refinement tool of initial solutions obtained either from a baseline MIQP with global linearization (Section~\ref{sec:preliminaries}) or historical schedules.

\begin{figure}[t!]
\centering
\includegraphics[width=\columnwidth]{figs/DFL_pipeline.pdf}
\caption{Decision-focused learning for UPHES day-ahead scheduling. Green arrows represent forward pass, while red arrows represent backward pass.}
\label{fig:methodology_framework}
\end{figure}

\begin{algorithm}[t!]
\caption{End-to-End Decision-Focused Learning}
\label{alg:dfl_scheduling}
\begin{algorithmic}[1]
\Statex \textbf{Input:} Training set $\mathcal{D} = \{(\boldsymbol{\lambda}_i^{\text{DA}}, \bar{\mathbf{x}}_i)\}_{i=1}^N$
\Statex \textbf{Output:} Trained parameters $\theta^*$
\vspace{1pt}
\Statex \hrule
\State Initialize $\theta^{(0)}$ 
\For{epoch $e = 1$ to $E$}
    \For{each sample $(\boldsymbol{\lambda}_i^{\text{DA}}, \bar{\mathbf{x}}_i)$ in $\mathcal{D}$}
        \State $\mathbf{w}_i^{(0)} \gets \exp(\mathcal{N}_{\theta^{(e)}}([\boldsymbol{\lambda}_i^{\text{DA}}, \bar{\mathbf{x}}_i]))$, $\hat{\mathbf{x}}_i^{(0)} \gets \bar{\mathbf{x}}_i$ 
        \For{$k = 0$ to $K-1$}
            \State $\mathbf{w}_i^{(k)} \gets \gamma^k \mathbf{w}_i^{(0)}$
            \State Linearize $f^{\text{UPC}}$ and $f^{\text{vol}}$ at $\hat{\mathbf{x}}_i^{(k)}$ to obtain $\boldsymbol{\xi}_i^{(k)}$
            \State Solve QP~\eqref{eq:opt_objective_penalty} with $\boldsymbol{\xi}_i^{(k)}$, $\mathbf{w}_i^{(k)}$ to get $\hat{\mathbf{x}}_i^{(k+1)}$
        \EndFor
        \State $\tilde{\mathbf{x}}_i \gets \textsc{Simulate}(\hat{\mathbf{x}}_i^{(K)})$ using Algorithm~\ref{alg:simulation}
        \State $\mathcal{L}_i \gets -\Pi_i(\tilde{\mathbf{x}}_i)$ via~\eqref{eq:expost_profit}
    \EndFor
    \State $\mathcal{L}(\theta^{(e)}) \gets \frac{1}{N}\sum_{i=1}^N \mathcal{L}_i$
    \State Update parameters: $\theta^{(e+1)} \gets \theta^{(e)} - \eta \nabla_\theta \mathcal{L}(\theta^{(e)})$
\EndFor
\State \Return $\theta^*$
\end{algorithmic}
\end{algorithm}

\subsection{Neural Penalty Prediction}
\label{subsec:penalty_prediction}

The neural penalty predictor estimates time-varying regularization weights that guide the linearization of noisy initial solutions. Given the day-ahead price vector $\boldsymbol{\lambda}^{\text{DA}} = [\lambda_1^{\text{DA}}, \ldots, \lambda_{24}^{\text{DA}}]^\top$ and initial noisy trajectories $\bar{\mathbf{x}} = \{\bar{p}_t, \bar{q}_t, \bar{h}_t\}_{t \in \mathcal{T}}$ from, e.g., a baseline MIQP solver, the network constructs a feature matrix $\mathbf{X} \in \mathbb{R}^{24 \times 4}$ where each row contains $[\lambda_t^{\text{DA}}, \bar{p}_t, \bar{q}_t, \bar{h}_t]$. This representation captures both market signals and the physical states of the reservoir. 

An LSTM network processes the sequential input to capture temporal dependencies in pricing patterns and operating schedules. The recurrent architecture recognizes multi-hour pumping and generation cycles typical of energy storage dispatch. The final hidden state at each time step passes through a linear projection to produce log-domain penalty weights

\begin{equation}
[\log \mathbf{w}_p, \log \mathbf{w}_q, \log \mathbf{w}_h] = \text{Linear}(\text{LSTM}(\mathbf{X})),
\label{eq:network_output_vector_explicit}
\end{equation}

which are exponentiated and bounded to prescribed intervals to ensure numerical stability and prevent penalty weights from dominating the objective. Log-domain parameterization guarantees strictly positive penalty weights while accommodating the wide dynamic range across different operating conditions.

\subsection{Local Linearization}
\label{subsec:local_linearization}

The recursive linearization process constructs first-order Taylor approximations of the nonlinear UPC and volume-head relationships around the current operating point. At iteration $k$ within the inner loop of Algorithm~\ref{alg:dfl_scheduling}, the operating point $\hat{\mathbf{x}}^{(k)} = \{\hat{p}_t^{(k)}, \hat{q}_t^{(k)}, \hat{h}_t^{(k)}\}_{t \in \mathcal{T}}$ defines the expansion point for local linear approximations. 

Following this approximation, the UPC mapping~\eqref{eq:upc-poly} becomes
\begin{equation}
q_t \approx \begin{bmatrix} p_t & h_t & 1 \end{bmatrix} \boldsymbol{\xi}_{q,t}^{(k)},
\label{eq:taylor_flow}
\end{equation}
for each time step $t$ and active mode $m$ where
\begin{equation}
\boldsymbol{\xi}_{q,t}^{(k)} = \begin{bmatrix}
\left. \frac{\partial f_m^{\text{UPC}}}{\partial p} \right|_{(\hat{p}_t^{(k)}, \hat{h}_t^{(k)})} \\[4pt]
\left. \frac{\partial f_m^{\text{UPC}}}{\partial h} \right|_{(\hat{p}_t^{(k)}, \hat{h}_t^{(k)})} \\[4pt]
f_m^{\text{UPC}}(\hat{p}_t^{(k)}, \hat{h}_t^{(k)}) - \nabla f_m^{\text{UPC}}(\hat{p}_t^{(k)}, \hat{h}_t^{(k)}) \begin{bmatrix} \hat{p}_t^{(k)} \\ \hat{h}_t^{(k)} \end{bmatrix}
\end{bmatrix}.
\label{eq:xi_flow}
\end{equation}
Similarly, the volume-head coupling~\eqref{eq:volume_head} is linearized as
\begin{equation}
v_{\text{low},t} \approx \begin{bmatrix} h_t & 1 \end{bmatrix} \boldsymbol{\xi}_{v,t}^{(k)},
\label{eq:taylor_volume}
\end{equation}
with coefficient vector $\boldsymbol{\xi}_{v,t}^{(k)} \in \mathbb{R}^2$ defined by the first-order Taylor expansion of $f_{\text{low}}^{\text{vol}}$ around $\hat{h}_t^{(k)}$:
\begin{equation}
\boldsymbol{\xi}_{v,t}^{(k)} = \begin{bmatrix}
\left. \frac{d f_{\text{low}}^{\text{vol}}}{d h} \right|_{\hat{h}_t^{(k)}} \\[4pt]
f_{\text{low}}^{\text{vol}}(\hat{h}_t^{(k)}) - \left. \frac{d f_{\text{low}}^{\text{vol}}}{d h} \right|_{\hat{h}_t^{(k)}} \hat{h}_t^{(k)}
\end{bmatrix}.
\label{eq:xi_volume}
\end{equation}

All partial derivatives are computed via automatic differentiation, ensuring accurate gradient information of the polynomial UPC~\eqref{eq:upc-poly} and cubic volume-head functions. When $|\hat{p}_t^{(k)}| < \epsilon$, indicating idle mode, $\boldsymbol{\xi}_{q,t}^{(k)}=0$ enforces $q_t = 0$. 

\subsection{Differentiable Optimization}
\label{subsec:diff_optimizer}

Using the linearization coefficients computed in the previous subsection, the optimization layer solves a convex quadratic program at each iteration of the inner loop. The QP reformulates the original UPHES scheduling problem~\eqref{eq:uphes_scheduling_problem} by incorporating adaptive regularization penalties while replacing nonlinear constraints with their local linearizations.
The modified objective function becomes
\begin{align}
\max_{p,q,h} \quad 
& \Delta t \sum_{t \in \mathcal{T}} \left( \lambda_t^{\text{DA}} p_t - C_{\text{op}} p_t^2 \right) - \sum_{t \in \mathcal{T}} \big[ w_{p,t} (p_t - \hat{p}_t^{(k)})^2 \notag \\
& \quad + w_{h,t} (h_t - \hat{h}_t^{(k)})^2 + w_{q,t} (q_t - \hat{q}_t^{(k)})^2 \big], \label{eq:opt_objective_penalty} 
\end{align}
where the penalty terms regularize deviations from the noisy initial trajectory $\hat{\mathbf{x}}^{(k)}$, creating an adaptive trust region that ensures linearization validity.

The nonlinear UPC~\eqref{eq:upc_constraints} and volume-head~\eqref{eq:volume_dynamics} relationships are replaced by their local linearizations~\eqref{eq:taylor_flow} and~\eqref{eq:taylor_volume}. Mode constraints prevent transitions during linearization
\begin{align}
p_t \in \begin{cases}
[p_{\min}^T(h_t), p_{\max}^T(h_t)] & \text{if } \hat{p}_t^{(k)} > 0,\\
[p_{\min}^P(h_t), p_{\max}^P(h_t)] & \text{if } \hat{p}_t^{(k)} < 0,\\
\{0\} & \text{if } \hat{p}_t^{(k)} = 0,
\end{cases} \\
    q_t = 0 \quad \text{ if } \hat{p}_t^{(k)} = 0. \label{eq:opt_mode}
\end{align}
All remaining constraints~\eqref{eq:head_bounds}--\eqref{eq:Target} remain unchanged. This mode-locking strategy eliminates integer variables, transforming the MINLP into a tractable convex QP while maintaining the discrete structure inherited from the initial solution.

The convex QP is implemented using CVXPY with CVXPYLayers~\cite{cvxpylayers2019} to enable gradient computation via implicit differentiation. The implicit function theorem provides gradient computation for optimal primal variables with respect to penalty weights without differentiating through iterative solver steps.

After solving, the updated trajectories $\hat{\mathbf{x}}^{(k+1)} = \{\hat{p}_t^{(k+1)}, \hat{q}_t^{(k+1)}, \hat{h}_t^{(k+1)}\}_{t \in \mathcal{T}}$ replace the previous iterate for the next inner loop iteration. Penalty weights are scaled by growth factor $\gamma^k$ at each iteration, progressively tightening constraints to encourage convergence through an annealing effect where early iterations explore broader regions while later iterations refine solutions locally.

\subsection{Simulation Layer and Ex-Post Profit Evaluation}
\label{subsec:simulation}

The simulation layer validates physical feasibility by enforcing the original nonlinear UPC model~\eqref{eq:upc-poly} and volume-head relationship~\eqref{eq:head_vol_relation}. Given the optimized power trajectory $\hat{\mathbf{x}}^{(K)} = \{\hat{p}_t^{(K)}\}_{t \in \mathcal{T}}$ from the final iteration, the simulator computes actual flow rates $\tilde{\mathbf{x}} = \{\tilde{p}_t, \tilde{q}_t, \tilde{h}_t, \tilde{v}_{\text{low},t}\}_{t \in \mathcal{T}}$ by evaluating the polynomial UPC at realized head values, which are updated sequentially according to reservoir mass balance
\begin{align}
\tilde{v}_{\text{low},0} &= v_{\text{low}}^{\text{init}} \label{eq:sim_v0}\\
\tilde{v}_{\text{low},t} &= \tilde{v}_{\text{low},t-1} + \tilde{q}_t \cdot \Delta t, \quad \forall t \in \mathcal{T} \label{eq:sim_vt}\\
\tilde{h}_t &= f_{\text{low}}^{\text{vol},-1}(\tilde{v}_{\text{low},t}), \quad \forall t \in \mathcal{T}. \label{eq:sim_ht}
\end{align}
Power outputs are clamped to head-dependent feasibility bounds to prevent UPC envelope violations. If predicted volume exceeds upper reservoir capacity or falls below the minimum admissible level, the unit is forced into idle mode. This clamping and mode switching logic ensures physical realizability even when the optimization layer produces infeasible suggestions due to approximation errors. Algorithm~\ref{alg:simulation} details this procedure.

\begin{algorithm}[t!]
\caption{Physical Simulation Layer}
\label{alg:simulation}
\begin{algorithmic}[1]
\Statex \textbf{Input:} Optimized power schedule $\hat{\mathbf{p}} = \{\hat{p}_t\}_{t \in \mathcal{T}}$
\Statex \textbf{Output:} Feasible simulated trajectories $\{\tilde{\mathbf{p}}, \tilde{\mathbf{q}}, \tilde{\mathbf{h}}, \tilde{\mathbf{v}}\}$
\vspace{1pt}
\Statex \hrule
\State Initialize: $\tilde{v}_{\text{low}} \gets v_{\text{low}}^{\text{init}}$, $\tilde{h} \gets h^{\text{init}}$
\For{$t = 1$ to $24$}
    \State Determine mode from $\hat{p}_t$ (idle, turbine, or pump)
    \State Clamp $\tilde{p}_t$ to $[p_{\min}^m(\tilde{h}_t), p_{\max}^m(\tilde{h}_t)]$
    \State Evaluate flow: $\tilde{q}_t \gets f_m^{\text{UPC}}(\tilde{p}_t, \tilde{h}_t)$
    \State Update volume: $\tilde{v}_{\text{low},t} \gets \tilde{v}_{\text{low},t-1} + \tilde{q}_t \cdot \Delta t$
    \If{$\tilde{v}_{\text{low},t} > v_{\max}$ or $\tilde{v}_{\text{low},t} < v_{\min}$}
        \State Force idle: $\tilde{p}_t \gets 0$, $\tilde{q}_t \gets 0$, $\tilde{v}_{\text{low},t} \gets \tilde{v}_{\text{low},t-1}$
    \EndIf
    \State Update head: $\tilde{h}_{t+1} \gets f_{\text{low}}^{\text{vol},-1}(\tilde{v}_{\text{low},t})$
\EndFor
\State \Return $\{\tilde{\mathbf{p}}, \tilde{\mathbf{q}}, \tilde{\mathbf{h}}, \tilde{\mathbf{v}}\}$
\end{algorithmic}
\end{algorithm}

The ex-post profit metric quantifies economic performance by accounting for revenue, operational costs, and penalties for constraint violations
\begin{align}
\Pi= \Delta t \sum_{t \in \mathcal{T}} \left( \lambda_t^{\text{DA}} \tilde{p}_t - C_{\text{op}} \tilde{p}_t^2 \right) - \text{SI} - \text{Vol},
\label{eq:expost_profit}
\end{align}
where the system imbalance penalty $\text{SI}$ accounts for deviations between optimized and simulated power outputs through asymmetric pricing: energy production shortages are penalized at twice the day-ahead price while surpluses are compensated only at half the day-ahead price, reflecting balancing market mechanisms. The volume penalty $\text{Vol}$ converts any excess water remaining beyond the target reservoir \eqref{eq:Target} level into equivalent energy value using system efficiency and hydraulic head parameters, priced at the median day-ahead rate to represent opportunity cost. No penalty is assessed when the terminal volume constraint is satisfied.

\subsection{Training Procedure}
\label{subsec:train_procedure}

Our goal is to train the neural penalty predictor offline using noisy MIQP solutions generated by perturbing baseline piecewise SOS2 optimization results. Perturbations are applied to power schedule variables $\{p_t\}_{t \in \mathcal{T}}$ while preserving operational modes $\{m_t \in \{I, T, P\}\}_{t \in \mathcal{T}}$ from the baseline solution. Perturbation magnitudes are uniformly sampled from intervals proportional to head-dependent capacity ranges $[p_{\min}^m(h_t), p_{\max}^m(h_t)]$, spanning noise levels from 10\% to 80\% as well as random sampling. Physically consistent flow rates $\{q_t\}_{t \in \mathcal{T}}$ are recovered by evaluating the polynomial UPC model~\eqref{eq:upc-poly} at the perturbed power values, and head trajectories $\{h_t\}_{t \in \mathcal{T}}$ are obtained by integrating volume dynamics~\eqref{eq:volume_dynamics} with the volume-head relationship~\eqref{eq:volume_head}. This generates training pairs $\{(\boldsymbol{\lambda}_i^{\text{DA}}, \bar{\mathbf{x}}_i) \mapsto \Pi_i\}_{i=1}^N$ where $\bar{\mathbf{x}}_i = \{\bar{p}_{i,t}, \bar{q}_{i,t}, \bar{h}_{i,t}\}_{t \in \mathcal{T}}$.

The loss function maximizes ex-post profit~\eqref{eq:expost_profit} obtained from the simulation layer after iterative linearization
\begin{equation}
\mathcal{L}(\theta) = -\frac{1}{N} \sum_{i=1}^N \Pi_i(\theta),
\label{eq:loss}
\end{equation}
where $\theta$ denotes network parameters. Training uses the Adam optimizer with plateau-based learning rate scheduling, early stopping, and gradient clipping.

\section{Results and Discussion}
\label{sec:results}

This section validates the proposed DFL framework through experiments on a realistic UPHES site in Belgium using actual electricity price scenarios presented in Section~\ref{sec:data}. Performance evaluation compares DFL methods against MIQP benchmarks from Section~\ref{sec:preliminaries}, and ablation studies are conducted to isolate architectural component contributions. Ex-post profit serves as the primary performance metric, computed by simulating optimized schedules under original system dynamics with realistic market penalties for constraint violations. 

\subsection{Data Description}
\label{sec:data}

The case study considers a 24-hour day-ahead scheduling problem with hourly time steps. As illustrated in Fig.~\ref{fig:uphes_system}, the UPHES plant features hydraulic head varying between $h_{\min} = 50$~m and $h_{\max} = 99$~m. As introduced in Section~\ref{sec:volume_constraints}, the upper reservoir exhibits a frustum geometry with an average slope $m=1.8$, while the lower reservoir consists of $n=100$ interconnected mine pits. Both reservoirs have maximum storage capacities of 588,000~m$^3$ and are initialized at 50\% capacity, with the terminal constraint enforcing $v_{\text{low}}^{\text{target}} = 294,000$~m$^3$ to ensure cyclical operation.
Day-ahead electricity prices $\lambda_t^{\text{DA}}$ are based on representative price scenarios in 2024, provided by Belgian transmission system operator Elia~\cite{EliaDayAheadReferencePrice}. Using k-medoids clustering, we identify 19 daily price profiles. The operational cost coefficient is set to $C_{\text{op}} = 0.4$ EUR/MW$^2$. The optimization is implemented using CVXPYlayers 0.1.6 with Gurobi 12.0.2 on an Intel Core Ultra 9 275HX Processor with 64 GB RAM running Windows 11. The MIP gap tolerance is set to 1\% with a time limit of 3600~s.

\subsection{Benchmark Approximation Quality}

Before evaluating the DFL framework, we characterize the approximation accuracy of the baseline MIQP methods introduced in Section~\ref{sec:preliminaries}. Table~\ref{tab:linearization_errors} shows approximation fidelity using four metrics evaluated on 480 operating points sampled from preliminary MIQP solutions across the 20 representative price scenarios. The metrics compare linearization predictions against evaluations using the true polynomial UPC~\eqref{eq:upc-poly} and cubic volume-head relationships~\eqref{eq:v_up_cubic},~\eqref{eq:v_low_cubic}.

\begin{table}[t!]
\centering
\caption{Benchmark Approximation Error Analysis}
\label{tab:linearization_errors}
\small
\begin{tabular}{@{}cccccc@{}}
\toprule
Function & Method & Mean & Max & MAPE & $R^2$ \\
 & & (\%) & (\%) & (\%) & \\
\midrule
\multirow{2}{*}{$f_m^{\text{UPC}}$} 
 & Piecewise Bilinear & 0.16 & 4.09 & 0.20 & 1.0000 \\
 & Global Linear & 4.36 & 18.31 & 5.78 & 0.9913 \\
\midrule
\multirow{2}{*}{$f^{\text{vol}}$} 
 & Piecewise Bilinear & 0.16 & 1.79 & 0.17 & 0.9999 \\
 & Global Linear & 15.98 & 18.08 & 16.08 & 0.4772 \\
\bottomrule
\end{tabular}
\end{table}

The piecewise bilinear approximation achieves MAPE below 0.21\% with near-perfect correlation ($R^2 \geq 0.9999$), while global linearization exhibits MAPE of 5.78\% for UPC and 16.08\% for volume-head mappings. The observed degradation in volume-head approximation accuracy stems from cubic geometry that cannot be captured by single-plane fits.

\subsection{Performance of Proposed Methods and Benchmarks}

\begin{figure*}[t!]
    \centering
    \includegraphics[width=\textwidth]{figs/profit_density_main_contribution.pdf} 
    \caption{Density distribution of ex-post profit comparing DFL methods against MIQP benchmarks across 19 representative price scenarios. Fig.~\ref{fig:penalty_tradeoff_ablation}(a) shows GL-based methods, Fig.~\ref{fig:penalty_tradeoff_ablation}(b) shows PW-based methods. Vertical dashed lines indicate mean values for each methodology. Densities estimated via kernel density estimation with Gaussian kernel (Scott's rule bandwidth, 0.5 scaling factor).}

    \label{fig:profit_density_main}
\end{figure*}

The performance of the proposed DFL framework is evaluated against MIQP benchmarks. Four methodologies are compared: piecewise-based DFL (PW-based DFL) employs piecewise bilinearization within the recursive linearization, global linearization-based DFL (GL-based DFL) utilizes global linear approximations, while benchmark methods MIQP-PW and MIQP-GL represent piecewise SOS2 bilinear and global linear formulations solved through mixed-integer optimization.

Training data for the DFL framework are constructed following the methodology described in Section~\ref{subsec:train_procedure}. Noise-injected MIQP baselines exhibit monotonic performance degradation, with MIQP-PW declining from 3837 EUR at 10\% noise to 3365 EUR under random sampling (12.3\% reduction) and MIQP-GL deteriorating from 3624 EUR to 2810 EUR (22.5\% reduction). These perturbed schedules serve as warm-start solutions, enabling the neural penalty predictor to learn linearization strategies that achieve near-optimal performance despite initialization corruption. All DFL methods reported in subsequent tables are trained using the random sampling noise level, representing the most challenging initialization scenario.

Table~\ref{tab:main_results} demonstrates that the DFL framework effectively navigates the trade-off between solution quality and computation time: as a refinement tool initialized with piecewise MIQP, it improves ex-post profit by 1.1\% over its baseline; alternatively, as a real-time scheduler initialized with linear approximations, it achieves a 300-fold speedup (3.87~s vs 1205.79~s) with 8.9\% profit improvement, while maintaining profitability within 3.6\% of the piecewise MIQP benchmark.

\begin{table}[t!]
\centering
\caption{Performance of DFL Methods and MIQP Benchmarks}
\label{tab:main_results}
\small
\begin{tabular}{lcc}
\toprule
\textbf{Method} & \textbf{Ex-post Profit (€)} & \textbf{Time (s)} \\
\midrule
MIQP-GL & $3409 \pm 1433$ & 2.69 \\
\textbf{DFL} (GL-based) & $\textbf{3712} \pm 1509$ & 2.69 + \textbf{1.18} \\
MIQP-PW & $3849 \pm 1545$ & 1205.79 \\
\textbf{DFL} (PW-based) & $\textbf{3890} \pm 1547$ & 1205.79 + \textbf{1.24} \\
\bottomrule
\end{tabular}
\end{table}

Fig.~\ref{fig:profit_density_main} shows the density distributions of ex-post profit across scenarios. For GL-based methods (Fig.~\ref{fig:penalty_tradeoff_ablation}(a)), DFL exhibits rightward-shifted distribution in the 3500--4000 EUR range compared to MIQP-GL. For PW-based methods (Fig.~\ref{fig:penalty_tradeoff_ablation}(b)), PW-based DFL and MIQP-PW display nearly overlapping distributions with DFL exhibiting slightly higher density in the primary region near 3900 EUR. The vertical dashed lines confirm the ranking PW-based DFL > MIQP-PW > GL-based DFL > MIQP-GL, where DFL improves solutions within each approximation framework through learned penalty prediction. The close proximity between piecewise methods indicates that DFL achieves near-optimal performance while substantially reducing computational time.

\begin{figure}[t!]
    \centering
    \includegraphics[width=\columnwidth]{figs/noise_robustness_ablation_study.pdf}
    \caption{Noise robustness comparison across perturbation levels. DFL methods (solid lines) maintain near-constant performance regardless of initialization quality.  MIQP benchmarks (dashed lines) are shown for reference.}
    \label{fig:noise_robustness}
\end{figure}

Fig.~\ref{fig:noise_robustness} demonstrates the noise-invariant performance of DFL, which maintains stable ex-post profits across 10\%--80\% perturbation levels (3870--3890 EUR for PW-based; 3710--3735 EUR for GL-based). In contrast, MIQP baselines exhibit a 12.3\%--22.5\% degradation under comparable conditions. This robustness confirms that the neural penalty predictor effectively compensates for initialization corruption, recovering near-optimal solutions from noisy initial guesses by guiding continuous optimization toward high-profitability regions.

\subsection{Ablation Study: Impact of Neural Network Components}

\begin{figure*}[t!]
    \centering
    \includegraphics[width=\textwidth]{figs/profit_vs_penalties_ablation.pdf}
    \caption{Penalty-profit trade-offs for DFL ablation variants. (a)~Ex-post profit versus volume penalty. (b)~Ex-post profit versus system imbalance penalty. Complete DFL (PW-based) and DFL (PW-no-Rec) achieve superior profit-penalty balance through learned constraint management.}
    \label{fig:penalty_tradeoff_ablation}
\end{figure*}

To isolate contributions of individual architectural components within the DFL (PW-based) framework, two ablated variants are evaluated: DFL (PW-no-Rec) eliminates recursive linearization, executing a single optimization pass, while DFL (PW-no-NN) removes the neural penalty predictor entirely, substituting a fixed heuristic penalty schedule.

Table~\ref{tab:ablation_results} shows these contributions, where all DFL variants are trained with random sampling initialization. Complete DFL (PW-based) achieves 3890 EUR mean profit in 1.24~s. Removing recursive linearization (DFL-no-Rec: 3851 EUR, 0.46~s) causes 1.0\% degradation, while eliminating the neural network (DFL-no-NN: 3780 EUR, 1.31~s) reduces performance by 2.8\%. The neural network component contributes 22.9\% of the total profit gain over MIQP-GL baseline, while recursive linearization adds an additional 8.1\%, establishing that the neural penalty predictor provides the dominant contribution to noise-robust optimization. Fig.~\ref{fig:noise_robustness} shows this conclusion is consistent across all perturbation levels.

\begin{table}[t!]
\centering
\caption{Ablation Study Results for DFL Components}
\label{tab:ablation_results}
\small
\begin{tabular}{lcc}
\toprule
\textbf{Method} & \textbf{Ex-post Profit (€)} & \textbf{Time (s)} \\
\midrule
DFL (PW-based) & $3890 \pm 1547$ & 1.24 \\
DFL (PW-no-NN) & $3780 \pm 1530$ & 1.31 \\
DFL (PW-no-Rec) & $3851 \pm 1532$ & 0.46 \\
\midrule
MIQP-PW (Benchmark) & $3849 \pm 1545$ & 1205.79 \\
\bottomrule
\end{tabular}
\end{table}

Fig.~\ref{fig:penalty_tradeoff_ablation}(a) reveals the adaptive constraint management enabled by learned penalty prediction. Complete DFL (PW-based) and DFL (PW-no-Rec) accept learned volume penalties averaging 250--280 EUR while achieving ex-post profits exceeding 3860 EUR. This results from the learnable trust region design combined with ex-post profit training, enabling the framework to discover that accepting calculable volume violations (priced at the median daily price) during high-revenue periods yields higher net profit. Conversely, DFL (PW-no-NN) incurs substantially higher volume penalties (averaging 515 EUR) and positive SI penalties (5--18 EUR) with lower profits around 3760--3810 EUR, indicating that fixed penalties produce constraint-violating solutions, producing schedules that violate both volume targets and power trajectory feasibility during physical simulation.

Fig.~\ref{fig:penalty_tradeoff_ablation}(b) demonstrates that DFL learns to exploit asymmetric market pricing, where underproduction incurs penalties at twice the day-ahead price, while overproduction receives compensation at half price. Complete DFL (PW-based) and DFL (PW-no-Rec) achieve negative SI penalty (averaging $-20$ EUR) by biasing schedules toward slight overproduction through learned penalty weights, accepting modest half-price revenue reductions to avoid costly double-price underproduction penalties while maintaining ex-post profits exceeding 3860 EUR. In contrast, DFL (PW-no-NN) incurs positive SI penalties of 5--18 EUR with lower profits around 3760--3810 EUR, indicating that fixed heuristic penalties fail to encode this asymmetric cost structure and result in schedules with underproduction. 

These ablation findings establish that learned penalty prediction constitutes the primary mechanism enabling noise-robust optimization and learned profit maximization, while recursive linearization provides incremental value in the profit. For time-critical applications requiring sub-second response, DFL (PW-no-Rec) presents an attractive alternative, retaining the majority of DFL's capabilities through neural penalty prediction alone.

\section{Conclusion}
\label{sec:conclusion}

This paper proposed a decision-focused learning framework for UPHES day-ahead scheduling that resolves the trade-off between approximation fidelity and computational speed inherent in mixed-integer nonlinear formulations. By integrating neural penalty prediction with recursive local linearization, the framework effectively navigates the trade-off between accuracy and computation. As a refinement tool, it improves profit by 1.1\% over piecewise MIQP baselines; alternatively, as a real-time scheduler initialized with global linear MIQP solutions, it achieves a 300-fold speedup while maintaining profitability within 3.6\% of the piecewise MIQP baselines.

Future research directions include direct end-to-end learning architectures that eliminate warm-start initialization requirements through physics-informed neural networks, and extension to intraday market participation, where shorter decision horizons and sequential market clearing mechanisms introduce distinct technical challenges, warranting dedicated investigation.

% ========================================
% BIBLIOGRAPHY
% ========================================
\bibliographystyle{IEEEtran}
\bibliography{bibtex/bib/IEEEabrv,reference}

\end{document}